%%%%%%%%%%%%%%%%%%%%%%%%%%%%%%%%%%%%%%%%%%%%%%%%%%%%%%%%%%%%%%%%%%%%%%
% Lớp 10 | Phương pháp tọa độ trong mặt phẳng (𝑂𝑥𝑦) | Tích vô hướng (theo tọa độ)
% Dạng 0H9?2-6: Cực trị và chứng minh bất đẳng thức
%%%%%%%%%%%%%%%%%%%%%%%%%%%%%%%%%%%%%%%%%%%%%%%%%%%%%%%%%%%%%%%%%%%%%%

% nguồn: TuLuan_Lop10_Dot2 (nhập 14/11/2026)
\begin{bt}%[0H9H2-6] % KIEM_TRA%[ma:ac95eacc]
	Trong mặt phẳng toạ độ $Oxy$, cho hai điểm $A(1; 2)$, $B(-4 ; 3)$. Gọi $M(t ; 0)$ là một điểm thuộc trục hoành.
	\begin{listEX}[2]
		\item Tính $\overrightarrow{A M} \cdot \overrightarrow{B M}$ theo $t$.\\
		\item Tìm $t$ để $\widehat{A M B}=90^{\circ}$.
	\end{listEX}
	\loigiai{
		\begin{enumerate}
			\item Tính $\overrightarrow{A M} \cdot \overrightarrow{B M}$ theo $t$.\\
			Ta có $\overrightarrow{A M}=(t-1 ;-2) ; \overrightarrow{B M}=(t+4 ;-3)$.
			\[\overrightarrow{A M} \cdot \overrightarrow{B M}=(t-1) \cdot(t+4)+(-2) \cdot(-3)=t^{2}+3 t+2.\]
			\item Tìm $t$ để $\widehat{A M B}=90^{\circ}$.
			\[\widehat{A M B}=90^{\circ} 
			\Leftrightarrow \overrightarrow{A M} \perp \overrightarrow{B M} 
			\Leftrightarrow \overrightarrow{A M} \cdot \overrightarrow{B M}=0 
			\Leftrightarrow t^{2}+3 t+2=0 
			\Leftrightarrow \hoac{& t=-1 \\ & t=-2.}\]
	\end{enumerate}}
\end{bt}

% nguồn: TuLuan_Lop10_Dot2 (nhập 14/11/2026)
\begin{bt}%[0H9V2-6] % KIEM_TRA%[ma:4235b44d]
	Cho điểm $A(2;1)$. Lấy điểm $B$ nằm trên trục hoành có hoành độ không âm và điểm $C$ trên trục tung có tung độ dương sao cho tam giác $ABC$ vuông tại $A$. Tìm toạ độ $B, C$ để tam giác $ABC$ có diện tích lớn nhất.
	\loigiai{
		Gọi $B(b;0), C(0;c)$ với $b\geq 0$, $c>0$. 
		Suy ra $\overrightarrow{AB}=(b-2;-1),\overrightarrow{AC}=(-2;c-1)$.
		\\
		Theo giả thiết ta có tam giác $ABC$ vuông tại $A$ nên  $$\overrightarrow{AB}\cdot\overrightarrow{AC}=0\Leftrightarrow(b-2)(-2)-1\cdot (c-1) =0 \Leftrightarrow c=-2b+5.$$
		Ta có $S_{\triangle ABC}=\dfrac{1}{2}AB\cdot AC=\dfrac{1}{2}\sqrt{(b-2)^2+1\cdot}\sqrt{2^2+(c-1)^2}=(b-2)^2+1=b^2-4b+5$.
		\\
		Vì $c>0$ nên $-2b+5>0\Rightarrow 0\leq b<\dfrac{5}{2}$.
		\\
		Xét hàm số $y=x^2-4x+5$ với $0\leq x<\dfrac{5}{2}$.
		\\
		Bảng biến thiên
		\begin{center}
			\begin{tikzpicture}
				\tkzTabInit[nocadre=false,lgt=1.2,espcl=2.5,deltacl=0.6]
				{$x$ /0.96,$y$ /2}
				{$0$,$2$,$\dfrac{5}{2}$}
				\tkzTabVar{+/$5$, -/$1$,+/$\dfrac{5}{4}$}
			\end{tikzpicture}
		\end{center}
		Suy ra giá trị lớn nhất của hàm số $y=x^2-4x+5$ với $0\leq x<\dfrac{5}{2}$ là $y=5$ khi $x=0$. \\ 
		Do đó diện tích tam giác $ABC$ lớn nhất khi và chỉ khi $b=0$, suy ra $c=5$.
		\\
		Vậy $B(0;0)$, $C(0;5)$ là điểm cần tìm.
	}
\end{bt}

% nguồn: TuLuan_Lop10_Dot2 (nhập 14/11/2026)
\begin{bt}%[0H9C2-6]%[Đề kiểm tra HK1 môn Toán 10 trường THPT Quảng Xương 4]%[Phạm Duy Phương, 10EX-HK1-2223]%[ma:e4192512]
Trong mặt phẳng tọa độ $Oxy$ cho các điểm $A(3;-1)$, $B(2;1)$ và $M$ là điểm thuộc trục $Ox$ sao cho $\left|3\overrightarrow{MA}-2 \overrightarrow{MB}\right|$ đạt giá trị nhỏ nhất. Tính độ dài bán kính đường tròn nội tiếp tam giác $ABM$ (làm tròn đến ba chữ số thập phân).
\loigiai{
	$M$ là điểm thuộc trục $Ox$ nên $M(x;0)$ $\Rightarrow \heva{&\overrightarrow{MA}=\left(3-x;-1\right)\\&\overrightarrow{MB}=\left(2-x;1\right).}$\\
	$3\overrightarrow{MA}-2\overrightarrow{MB}=\left(5-x;-5\right) \Rightarrow \left|3\overrightarrow{MA}-2\overrightarrow{MB}\right|=\sqrt{(5-x)^2+25} \geq 5,\, \forall x$.\\
	Do đó $\left|3\overrightarrow{MA}-2 \overrightarrow{MB}\right|$ đạt giá trị nhỏ nhất là $5$ khi $5-x=0 \Leftrightarrow x=5$.\\
	Suy ra $M(5;0)$ $\Rightarrow \heva{&\overrightarrow{MA}=\left(-2;-1\right) \Rightarrow \left|\overrightarrow{MA}\right|=\sqrt{5}\\&\overrightarrow{MB}=\left(-3;1\right) \Rightarrow \left|\overrightarrow{MB}\right|=\sqrt{10}\\&\overrightarrow{AB}=\left(-1;2\right) \Rightarrow\left|\overrightarrow{AB}\right|=\sqrt{5}.}$\\
	$\triangle MAB$ có $MB^2=MA^2+AB^2$ nên vuông tại $A$ suy ra diện tích $\triangle MAB$ là $$S=\dfrac{1}{2}\cdot \sqrt{5}\cdot \sqrt{5}=\dfrac{5}{2}.$$
	Nửa chu vi của $\triangle MAB$ là $p=\dfrac{\sqrt{5}+\sqrt{10}+\sqrt{5}}{2}$.\\
	Gọi $r$ là độ dài bán kính đường tròn nội tiếp tam giác $ABM$ ta có 
	$$S=pr \Leftrightarrow r=\dfrac{S}{p}=\dfrac{\tfrac{5}{2}}{\tfrac{\sqrt{5}+\sqrt{10}+\sqrt{5}}{2}}=\dfrac{1}{2\sqrt{2}+1} \approx 0{,}261.$$
}
\end{bt}

% nguồn: TuLuan_Lop10_Dot2 (nhập 14/11/2026)
\begin{bt}%[0H9V2-6]%[ma:7f150766]
	Cho ba điểm $A(-1;4)$, $B(1;1)$, $C(3;-1)$. Tìm điểm $M$ thuộc trục hoành sao cho $\left|MA-MB\right|$ bé nhất.
\loigiai{
	Ta thấy $y_A\cdot y_B=4.1>0 \Rightarrow$ $A$, $B$ nằm cùng phía so với trục $Ox$. Ta có $|AM-BM| \leq AB$ nên $|AM-BM|_{\max}=AB$.\\
	Giá trị lớn nhất này đạt được khi $A$, $B$, $M$ thẳng hàng ($M$ nằm ngoài $AB$).\\
	Gọi $M(x;0) \in Ox \Rightarrow \overrightarrow{AM}=(x+1;-4), \overrightarrow{AB}=(2;-3)$.\\
	Ta có $\overrightarrow{AM}$ và $\overrightarrow{AB}$ cùng phương $\Leftrightarrow \dfrac{x+1}{2}=\dfrac{-4}{-3} \Leftrightarrow 3(x+1)=8 \Leftrightarrow x=\dfrac{5}{3}$ hay $M\left(\dfrac{8}{3}; 0\right)$.
}
\end{bt}

% nguồn: TuLuan_Lop10_Dot2 (nhập 14/11/2026)
\begin{bt}%[0H9V2-6]%[1261-1280 Sỹ Trường]%%[ma:562631ec]
	Trong mặt phẳng tọa độ $Oxy$, cho $M$ thuộc trục $Ox$ và $A(4; 1), B(1; 3)$. Tìm giá trị nhỏ nhất của biểu thức $T = MA + MB$.
	\loigiai
	{
		\immini{Ta có $A(4; 1), B(1; 3)$ nằm về cùng một phía đối với trục $Ox$.
			\\ Gọi $C$ đối xứng với $A$ qua trục $Ox$ suy ra $C(4; -1)$.
			\\ Với mọi điểm $M$ thuộc trục $Ox$ ta có $T= MA+MB = MC + MB \geq BC = 5$.}
		{\begin{tikzpicture}[>=stealth,scale=1, line join = round, line cap = round]
				\draw[->,line width = 1pt] (-1,0)--(0,0) node[below right]{$O$}--(5,0) node[below]{$x$};
				\draw[->,line width = 1pt] (0,-1.5) --(0,4) node[right]{$y$};
				\tikzset{label style/.style={font=\footnotesize}}
				\tkzDefPoints{4/1/A,1/3/B,4/-1/C,2/0/M}
				\tkzDrawSegments(M,A M,B M,C B,C)
				\tkzDrawSegments[dashed](A,C)
				\tkzDrawPoints(A,B,C,M)
				\tkzLabelPoints[left](B)
				\tkzLabelPoints[right](A,C)
				\tkzLabelPoints[below](M)
	\end{tikzpicture}}}
\end{bt}

% nguồn: De_AI_Toan10_cuoi_hoc_ki_i_1 (nhập 21/11/2026)
\begin{bt}%[0H9C2-6]%[ma:67a3ea5f]
	Trong mặt phẳng tọa độ $Oxy$, cho hai điểm $A(1; 2)$ và $B(5; 6)$. Xét điểm $M$ di động trên trục hoành $Ox$ sao cho tổng độ dài $MA + MB$ đạt giá trị nhỏ nhất. Khi đó, hãy tính tích vô hướng $\overrightarrow{MA} \cdot \overrightarrow{MB}$.
	\shortans{9}
	\loigiai{
		Bước 1: Vì $y_A = 2 > 0$ và $y_B = 6 > 0$ nên hai điểm $A, B$ nằm cùng phía đối với trục hoành $Ox$. \\
Bước 2: Gọi $B'(5; -6)$ là điểm đối xứng với $B$ qua trục $Ox$. Khi đó với mọi điểm $M \in Ox$, ta có $MB = MB'$ nên $MA + MB = MA + MB' \ge AB'$. \\
Bước 3: Dấu bằng xảy ra khi và chỉ khi $M$ là giao điểm của đoạn thẳng $AB'$ với trục $Ox$, tức là $A, M, B'$ thẳng hàng. Đặt $M(a; 0)$, ta có $\overrightarrow{AM} = (a - 1; -2)$ và $\overrightarrow{AB'} = (4; -8)$. \\
Bước 4: Hai vectơ $\overrightarrow{AM}$ và $\overrightarrow{AB'}$ cùng phương khi $\dfrac{a - 1}{4} = \dfrac{-2}{-8} \Leftrightarrow a - 1 = 1 \Leftrightarrow a = 2$. Suy ra $M(2; 0)$. \\
Bước 5: Ta có $\overrightarrow{MA} = (-1; 2)$ và $\overrightarrow{MB} = (3; 6)$. Do đó $\overrightarrow{MA} \cdot \overrightarrow{MB} = (-1) \cdot 3 + 2 \cdot 6 = 9$.
	}
\end{bt}

% nguồn: De_AI_Toan10_cuoi_hoc_ki_i_1 (nhập 21/11/2026)
\begin{bt}%[0H9C2-6]%[ma:a09d1389]
	Trong mặt phẳng tọa độ $Oxy$, cho điểm $I(1; 2)$ và điểm $M(x; y)$ thay đổi sao cho độ dài vectơ $\overrightarrow{IM}$ luôn bằng $2$. Tìm giá trị lớn nhất của biểu thức $P = 3x + 4y$.
	\shortans{21}
	\loigiai{
		Bước 1: Ta có $\overrightarrow{IM} = (x - 1; y - 2)$ và $|\overrightarrow{IM}| = 2$. \\
Bước 2: Xét vectơ $\overrightarrow{u} = (3; 4)$, ta có $|\overrightarrow{u}| = \sqrt{3^2 + 4^2} = 5$. \\
Bước 3: Tính tích vô hướng $\overrightarrow{u} \cdot \overrightarrow{IM} = 3(x - 1) + 4(y - 2) = 3x + 4y - 11 = P - 11$. \\
Bước 4: Theo bất đẳng thức về tích vô hướng, ta có $P - 11 = \overrightarrow{u} \cdot \overrightarrow{IM} \le |\overrightarrow{u}| \cdot |\overrightarrow{IM}| = 5 \cdot 2 = 10$, suy ra $P \le 21$. \\
Bước 5: Dấu bằng xảy ra khi $\overrightarrow{IM}$ cùng hướng với $\overrightarrow{u}$, tức là $\overrightarrow{IM} = \dfrac{2}{5}\overrightarrow{u} = \left(\dfrac{6}{5}; \dfrac{8}{5}\right)$, tương ứng $M\left(\dfrac{11}{5}; \dfrac{18}{5}\right)$. Vậy giá trị lớn nhất của $P$ là $21$.
	}
\end{bt}

% nguồn: 2026-11-21_193956_hk_ai_lop10 (nhập 21/11/2026)

% nguồn: 2026-11-21_193956_hk_ai_lop10 (nhập 21/11/2026)

% nguồn: 2026-11-21_200956_hk_ai_lop10 (nhập 21/11/2026)
\begin{bt}%[0H9V2-6]%[ma:a3a099cf]
% nguon: AI tự sinh (Ôn thi HK)
	Trong mặt phẳng $Oxy$, một robot di chuyển trên đường thẳng $\Delta\colon x-2y+4=0$. Tại mỗi vị trí $M$ trên $\Delta$, robot đo khoảng cách đến hai trạm phát sóng đặt tại $A(1;1)$ và $B(5;3)$. Tìm giá trị nhỏ nhất của tổng $MA^2+MB^2$ (làm tròn đến hàng phần mười).
	\shortans{13,6}
	\loigiai{
		Gọi $I$ là trung điểm $AB$: $I(3;2)$. Ta có $MA^2+MB^2=2MI^2+\dfrac{AB^2}{2}$.\\
		$AB^2=(5-1)^2+(3-1)^2=16+4=20$, nên $MA^2+MB^2=2MI^2+10$.\\
		Tổng nhỏ nhất khi $MI$ nhỏ nhất, tức $M$ là hình chiếu của $I$ lên $\Delta$.\\
		$MI_{\min}=d(I,\Delta)=\dfrac{|3-2\cdot 2+4|}{\sqrt{1+4}}=\dfrac{3}{\sqrt5}$.\\
		Vậy $\min(MA^2+MB^2)=2\cdot\dfrac{9}{5}+10=\dfrac{18}{5}+10=13{,}6$.
	}
	
\end{bt}

% nguồn: 2026-11-21_201230_hk_ai_lop10 (nhập 21/11/2026)
\begin{bt}%[0H9V2-6]%[ma:3c4d5e6f]
% nguon: AI tự sinh (Ôn thi HK)
Trong mặt phẳng $Oxy$, cho hai điểm $A(1;2)$, $B(4;-2)$ và điểm $M$ thay đổi trên đường thẳng $d\colon x-y+1=0$. Tìm giá trị nhỏ nhất của biểu thức $P=MA^2+MB^2$ (làm tròn đến hàng phần mười).
\shortans{24,8}
	\loigiai{
Gọi $I$ là trung điểm $AB$: $I\left(\dfrac{5}{2};0\right)$.\\
Theo công thức trung tuyến: $MA^2+MB^2=2MI^2+\dfrac{AB^2}{2}$.\\
$AB^2=(4-1)^2+(-2-2)^2=9+16=25\Rightarrow\dfrac{AB^2}{2}=12{,}5$.\\
Vậy $P$ nhỏ nhất khi $MI$ nhỏ nhất, tức $M$ là hình chiếu của $I$ lên $d$.\\
$MI_{\min}=d(I,d)=\dfrac{\left|\frac52-0+1\right|}{\sqrt{1+1}}=\dfrac{7/2}{\sqrt2}=\dfrac{7}{2\sqrt2}$.\\
$MI_{\min}^2=\dfrac{49}{8}=6{,}125$.\\
$P_{\min}=2\cdot6{,}125+12{,}5=12{,}25+12{,}5=24{,}75\approx24{,}8$.
}

\end{bt}

% nguồn: 2026-11-21_201230_hk_ai_lop10 (nhập 21/11/2026)
\begin{bt}%[0H9V2-6]%[ma:9c0d1e2f]
% nguon: AI tự sinh (Ôn thi HK)
Trong mặt phẳng $Oxy$, cho điểm $A(1;1)$ và đường thẳng $d\colon x+2y-4=0$. Điểm $M$ thay đổi trên $d$. Tìm giá trị nhỏ nhất của $MA$ (làm tròn đến hàng phần trăm).
\shortans{0,45}
	\loigiai{
$MA$ nhỏ nhất khi $M$ là hình chiếu vuông góc của $A$ lên $d$.\\
$MA_{\min}=d(A,d)=\dfrac{|1+2-4|}{\sqrt{1+4}}=\dfrac{1}{\sqrt5}=\dfrac{\sqrt5}{5}\approx0{,}447\approx0{,}45$.
}

\end{bt}

% nguồn: 2026-11-21_201230_hk_ai_lop10 (nhập 21/11/2026)
\begin{bt}%[0H9V2-6]%[ma:5a7e93bc]
% nguon: AI tự sinh (Ôn thi HK)
Trong mặt phẳng $Oxy$, cho hai điểm $A(1;3)$ và $B(4;-1)$. Điểm $M$ thay đổi trên đường thẳng $d\colon x+y-2=0$. Gọi $T=MA^2+MB^2$. Tìm giá trị nhỏ nhất của $T$ (làm tròn đến hàng phần mười).
\shortans{14,8}
	\loigiai{
Gọi $M(x;y)\in d\Rightarrow y=2-x$, tức $M(x;2-x)$.\\
Ta có $MA^2=(x-1)^2+(2-x-3)^2=(x-1)^2+(x+1)^2=2x^2+2$.\\
Và $MB^2=(x-4)^2+(2-x+1)^2=(x-4)^2+(3-x)^2=2x^2-14x+25$.\\
Suy ra $T=MA^2+MB^2=(2x^2+2)+(2x^2-14x+25)=4x^2-14x+27$.\\
Tam thức bậc hai $4x^2-14x+27$ đạt giá trị nhỏ nhất tại $x=\dfrac{14}{8}=\dfrac{7}{4}$.\\
Khi đó $T_{\min}=4\cdot\dfrac{49}{16}-14\cdot\dfrac{7}{4}+27=\dfrac{49}{4}-\dfrac{49}{2}+27=\dfrac{49-98+108}{4}=\dfrac{59}{4}=14{,}75\approx 14{,}8$.\\
Vậy giá trị nhỏ nhất của $T$ xấp xỉ $14{,}8$.
}

\end{bt}
